\documentclass[11pt,letterpaper]{article}
\usepackage[margin=1in]{geometry}
\usepackage[T1]{fontenc}
\usepackage{lmodern,amsmath,amssymb,microtype}
\usepackage[colorlinks=true,linkcolor=blue,urlcolor=blue]{hyperref}
\hypersetup{pdftitle={Verification of the Appendix C.2 Certificate},
 pdfsubject={Finite dual certificate for precommitted thresholds in the i.i.d. setting}}
\setlength{\parindent}{0pt}
\setlength{\parskip}{5pt}
\title{Verification of the Appendix C.2 Certificate}
\author{}
\date{}
\begin{document}
\maketitle
\vspace{-2em}

\section{Certificate and coefficient formulas}
The certificate consists of a finite collection $\mathcal Q$ of step rate
functions and weights $(\lambda_q)_{q\in\mathcal Q}$. Each
$\mathbf{q}=(q_0,\ldots,q_{N-1})$ represents
\[
 q(t)=q_j\quad\text{for }t\in[j/N,(j+1)/N),\qquad q(1)=q_{N-1}.
\]
The input checks establish
\[
 0\le q_0\le\cdots\le q_{N-1},\qquad
 \lambda_q\ge0,\qquad \sum_{q\in\mathcal Q}\lambda_q=1.
\]
All rates and weights are exact rational numbers, read from integer
numerators and positive integer denominators.

Recall the paper's definitions
\[
 S_q(t)=\exp\!\left(-\int_0^t q(s)\,ds\right),\qquad
 A_q(t)=\int_t^1 S_q(s)\,ds,\qquad
 a_j=\int_{j/N}^{(j+1)/N}S_q(t)\,dt.
\]
For a fixed $q$, the verifier uses
\begin{align*}
 S_q(0)&=1,\qquad
 S_q((j+1)/N)=S_q(j/N)e^{-q_j/N},\\
 a_j&=
 \begin{cases}
 S_q(j/N)\,(1-e^{-q_j/N})/q_j,&q_j>0,\\
 S_q(j/N)/N,&q_j=0,
 \end{cases}
 \qquad A_q(0)=\sum_{j=0}^{N-1}a_j,\\
 k_{q,j}&=S_q((j+1)/N)-q_j\sum_{\ell=j+1}^{N-1}a_\ell
                           -A_q(0)e^{-q_j}.
\end{align*}
The formula for $a_j$ when $q_j>0$ equals
$(S_q(j/N)-S_q((j+1)/N))/q_j$ by the recurrence for $S_q$.
Thus these are exactly the coefficient formulas of Appendix C.2.
The empty sum for $j=N-1$ is zero.

The numerical task is to verify the $N+1$ inequalities in Lemma C.11:
\[
 \sum_{q\in\mathcal Q}\lambda_q A_q(0)\le U,
 \qquad
 \sum_{q\in\mathcal Q}\lambda_q(1-k_{q,j})\le U
 \quad(0\le j<N).
\]

\newpage
\section{Rigorous interval evaluation}
Set $B=2^{160}$.
An integer pair $(\ell,u)$ represents $[\ell/B,u/B]$. Integer operations
are performed without overflow. Floor and ceiling have their usual
mathematical meanings, including for negative arguments.

\paragraph{Arithmetic operations.}
For input pairs $(\ell_1,u_1)$ and $(\ell_2,u_2)$, addition and subtraction
return
\[
 (\ell_1+\ell_2,u_1+u_2),\qquad
 (\ell_1-u_2,u_1-\ell_2),
\]
respectively. For multiplication, define
\begin{align*}
 P_{\min}&=\min\{\ell_1\ell_2,\ell_1u_2,u_1\ell_2,u_1u_2\},\\
 P_{\max}&=\max\{\ell_1\ell_2,\ell_1u_2,u_1\ell_2,u_1u_2\}.
\end{align*}
The returned pair is
\[
 \left(\left\lfloor P_{\min}/B\right\rfloor,
       \left\lceil P_{\max}/B\right\rceil\right).
\]
Multiplication by an exact nonnegative rational number multiplies the
two endpoints by that number, then rounds the lower endpoint down and
the upper endpoint up. For a negative multiplier, the endpoints are
first exchanged. Division by a positive rational is multiplication by
its exact reciprocal. The separate $q_j=0$ branch avoids division by zero.

These operations preserve containment. For multiplication, the extrema
of the product over the two intervals occur at their corners and equal
$P_{\min}/B^2$ and $P_{\max}/B^2$. The stated rounding can only enlarge
that range. The assertions for addition, subtraction, and rational
scaling follow directly from the order of their endpoints. This reasoning
does not require independent inputs: repeated occurrences of the same
exact quantity may widen an interval but cannot invalidate containment.

\paragraph{Exponential enclosure.}
Let $x\ge0$ be rational. If $x=0$, return $(B,B)$. If $x\ge160$,
return $(0,1)$; this is valid because
\[
 0<e^{-x}\le e^{-160}<2^{-160}=B^{-1}.
\]
For $0<x<160$, write $x=2^h z$ with an integer $h\ge0$ and $0<z\le1$.
Starting with $\ell_0=u_0=B$, compute for $m=1,\ldots,50$
\[
 \ell_m=\left\lfloor\frac{z\ell_{m-1}}m\right\rfloor,
 \qquad
 u_m=\left\lceil\frac{zu_{m-1}}m\right\rceil.
\]
Each rational expression inside floor or ceiling is evaluated exactly.
Induction, using multiplication by $z/m>0$, gives
\[
 \ell_m/B\le z^m/m!\le u_m/B\qquad(0\le m\le50).
\]

\newpage
The remainder after the degree-$50$ Taylor polynomial satisfies
\begin{align*}
 0\le e^z-\sum_{m=0}^{50}\frac{z^m}{m!}
 &=\sum_{m=51}^{\infty}\frac{z^m}{m!}\\
 &\le\frac1{51!}\sum_{r=0}^{\infty}\left(\frac1{52}\right)^r
 =\frac{52}{51\cdot51!}<\frac3{51!}.
\end{align*}
Here $z\le1$ and $(51+r)!\ge51!\,52^r$. Therefore the integers
\[
 L=\sum_{m=0}^{50}\ell_m,\qquad
 H=\sum_{m=0}^{50}u_m+\left\lceil\frac{3B}{51!}\right\rceil
\]
satisfy $0<L/B\le e^z\le H/B$. Taking reciprocals and rounding outward
yields
\[
 e^{-z}\in\left[
    \frac{\lfloor B^2/H\rfloor}{B},
    \frac{\lceil B^2/L\rceil}{B}\right].
\]
Starting from this nonnegative interval, square it $h$ times, rounding
outward after each step. On a nonnegative interval, squaring is
monotone, so the update for $(\ell,u)$ is
$(\lfloor\ell^2/B\rfloor,\lceil u^2/B\rceil)$.
Induction then proves containment of $(e^{-z})^{2^h}=e^{-x}$.

\paragraph{Coefficient evaluation.}
Apply these enclosures to the exact rational arguments $q_j/N$ and
$q_j$, and evaluate the formulas in Section~1 with the interval operations
above. An induction over the arithmetic operations shows that each
computed interval contains its exact counterpart. The verifier computes
the sum $\sum_{\ell=j+1}^{N-1}a_\ell$ by traversing $j$ in reverse order;
the current $a_j$ is added only \emph{after} evaluating $k_{q,j}$.
Reusing an enclosure for the same rational argument changes neither its
meaning nor its validity.

\section{Certification of the upper bound}
After multiplication by the exact weights and summation, write the
computed enclosures as
\[
 \sum_{q\in\mathcal Q}\lambda_q A_q(0)
       \in[\underline A,\overline A],\qquad
 \sum_{q\in\mathcal Q}\lambda_q k_{q,j}
       \in[\underline\kappa_j,\overline\kappa_j].
\]
Define
\[
 \widehat U=\max\left\{\overline A,
                      \max_{0\le j<N}(1-\underline\kappa_j)\right\}.
\]
Since the exact input checks establish $\sum_q\lambda_q=1$,
\[
 \sum_{q\in\mathcal Q}\lambda_q(1-k_{q,j})
 =1-\sum_{q\in\mathcal Q}\lambda_q k_{q,j}
 \le1-\underline\kappa_j\le\widehat U.
\]
The remaining average is at most $\overline A\le\widehat U$.
Consequently, if the exact comparison $\widehat U\le U$ passes, all
coefficient inequalities and all other hypotheses of Lemma C.11 hold.

\newpage
\section{Released data and verification result}
The file \texttt{data/upper\_1024.json.gz} contains $N=1024$,
$1014$ rate vectors, and their weights. A rate is its stored integer
divided by $10^8$, and a weight is its stored integer divided by
$10^{12}$. The verifier checks dimensions, these denominators,
rate nonnegativity and monotonicity, weight nonnegativity, and an exact
weight-numerator sum of $10^{12}$. It also checks that all rates are at
most $10^9$ and reports their actual maximum, which is $10^9$.
The auxiliary \texttt{cutoffs} field is retained from the original artifact
and is not used to evaluate or validate the certificate.

Running the released verifier on these data gives
\[
 \widehat U=
 \frac{1084385539732742468991572064431813102297740351518}
      {1461501637330902918203684832716283019655932542976}.
\]
The denominator is $B=2^{160}$. Exact integer comparisons establish
\[
 \widehat U\le
 \frac{741966695100}{10^{12}}
 =0.741966695100
 <U=0.741967.
\]
In particular, the verifier checks the first comparison by multiplying
the computed numerator by $10^{12}$ and comparing it with
$741966695100\,B$. It checks the second comparison as
$741966695100<741967000000$. Both comparisons are exact; the displayed
decimal bound is obtained by upward rounding, not by nearest rounding
or a floating-point feasibility tolerance.

The accompanying verification record gives the run environment,
the data and verifier hashes, and the computed bound. The coefficient
table records interval endpoints for all $1025$ coefficient averages.
The README identifies these files and the release's provenance and
checksum records.

All arithmetic affecting the certificate uses Python integers; rational
scaling uses exact integer division with directed rounding. No numerical
solver, floating-point exponential evaluation, or third-party Python
package is needed to verify the certificate. The release's README gives
the command for reproducing these checks.

The certificate therefore satisfies Lemma C.11 with $U=0.741967$.
Together with Lemma C.12, this yields the upper bound stated in Appendix
C.2. For this certificate the latter lemma's size condition is satisfied
when $n\ge1{,}000{,}000{,}001$. For every admissible threshold and every
$\eta>0$, Lemma C.12 supplies a normalized hard input with expected
reward at most $e^{1/n}U+\eta$; hence $U$ is the resulting asymptotic
upper bound. Verification establishes a feasible finite dual witness;
it does not require that the search procedure find an optimal one.

\end{document}
